\head{第二十版回校：第12节对第2--5节的推广}
\source{法文\texttt{EGA5\_francais.pdf}物理第66--69页；\\英文\texttt{EGAVaiello.pdf}物理第44--45页。另核法第16、18--20、23页的回指及正式\texttt{EGAIV.2.pdf}物理第142--143页（印刷第145--146页）}

\subsec{实际核对范围与段落对应}
本模块从法第66页中部“\emph{La proposition 2.1. se transpose sans changement}”开始，止于法第69页上部“\emph{Th\'eor\`eme 5.6. s'\'etend tel quel, de m\^eme 5.7}”；下一句命题6.1不属本模块。实际查看了法第66、67、68、69页及英第44、45页的整页图像，并由原PDF另渲染法第67、68页的第4、5节段落与第69页5.5段的放大图。法第67、68页起讫相同、改字和划销位置相同，是同一稿页的不同扫描取景；两份均核，正文只译一次。未发现可据以视作两个稿本的实质异文。

\begin{longtable}{@{}p{.17\textwidth}p{.39\textwidth}p{.34\textwidth}@{}}
\toprule 原页 & 原文内容 & 本版处理\\\midrule\endhead
法66；英44 & 2.1、2.2；独立一般超平面的归约；给Dieudonn\'e的插语 & 完整译出；恢复点闭包横线，保留插语及英译者疑问。\\
法66；英44--45 & 2.6同步编写、$\dim X-m$及负维数约定 & 保留作者的编写建议，不改成任意截面的维数定理。\\
法66；英45 & 2.3--2.5、2.7；2.8--2.11 & 逐项保留原序，明确2.7仍使用关联素循环$Z_i$。\\
法66末接67/68首；英45 & 2.12--2.18及第8节回指 & 跨页合为一次译文；手写“下文”译出。\\
法67/68；英45 & 3.1及引理3.1.1编写选择 & 恢复“除非……”关系及省略号；不补引理。\\
法67/68；英45 & 3.2、3.3的两个下界及闭集族 & 保留全部不等式；闭集族沿用$\sigma$。\\
法67/68；英45 & 第4节、4.2、4.3(b)(v)(vi)、4.4(b) & 正文用$m+1$；数学裁定与尚存字形歧义分开记录。\\
法67/68；英45 & 5.1；5.2中$T$、$G_T(1)$、$\phi^{(m)}_\xi$ & 核对载体及两次换底，划销层另记。\\
法67/68末接69首；英45 & 秩$m$、$m$个分量、基无关性待验、5.3及5.4 & 据法文改正英译$n$；保留待验提示。\\
法69上部；英45末 & 5.5两部分、$(S_m)$、5.6、5.7 & 恢复$F_s$；全部保留为作者的推广指示。\\
\bottomrule
\end{longtable}

旧版对应于\texttt{v\_dot\_middle.tex}的DOT-MIDDLE批03、04。Dot118是另一个待审核候选，不作为法文的替代底本。其“译者附校”没有混入本版原文正文；本报告以下论证标为译校论证，所补条件及引用均另行核定。候选全文及旧源由总工程保全，本报告末尾另逐字保存本范围旧注D17-M013至D17-M024。

\audit{V20-EXT-A01}{法文、英译与旧译的错误归属}{已分层订正，未把英译差错倒记为法稿错误}
法第66页的$\dim\bar x\leq m-1$在英文变成$\dim x\leq m-1$；正文恢复闭包。法第67/68页的\emph{\`A moins qu'il soit amusant\ldots}在英文成为“At least”，旧译因而作“至少”；现改为“除非觉得……”。法稿一般截面的局部自由模秩与分量数同为$m$，英文先写秩$n$再写$m$个分量；正文据法稿写$m$，不再把它描述为作者自身的秩数冲突。法第69页$F_s$的下标是小写，英文$F_S$须改正。

法第16页3.3与第67/68页的同一手写闭集族记号，在英文前处为$\sigma$、后处为$G$。采用$\sigma$是与被引定义统一的转写选择；这里没有新造闭集族，也没有把它等同于本节泛商模$G$。像维数$\dim f(W)$仍是原稿的简记，按前文解释为$\dim\overline{f(W)}$；它与逐点闭包维数条件有不同量词。

英文“[Le laius 5.1 Fr or Latin 67]”保存的是编者不确定的语言及稿页标记。法文可直接读为\emph{Le la\"ius 5.1. passe tel quel.}，正文译为“5.1中的那段议论照旧”。英文关于\emph{se transpose}的问号与\emph{Dieudonn\'e demerdetur}的语种疑问均收入相应脚注，没有消失在顺译中。

\subsec{划销和加笔层}
法第66页2.2的维数式前有被覆写的短串，无法可靠转写；式中的点闭包横线与$\leq m-1$可辨。证明中的旧词\emph{section}被划掉，现行措辞为一般的余维$m$线性子簇。2.3的维数界附近另有被覆写的多余字符；现行$\leq m-1$由该行与下一行连续读出。法第67/68页2.12的括注上方补入\emph{plus bas}，已译作“下文”。

法第67/68页5.2中，$T$后的\emph{est un pr\'esch\'ema}及下一行\emph{au-dessus de}被划销，取而代之的等式是$T=X\times_S\operatorname{Gr}^m$。等式后与纤维模式前还有重度覆写的旧短串及旧对象；无法可靠完整复原，不把猜测当正文。纤维模的现行式按$\mathcal O_{X_{k(\xi)}}(1)\otimes_{k(\xi)}G(\xi)$读，清楚保留其载体、扭转与泛商纤维。法第69页5.5首句把“推广到”的措辞移到下一行，并以手写$(S_m)$替换被划掉的条件；被覆盖缩写似为“CM”，但字形仍存疑，故只把确定可读的修订条件用于正文。这里未发现可独立恢复的一整段被删证明；原页及上述覆盖位置均保留供复核。

\audit{V20-EXT-A02}{\texorpdfstring{4.3(b)的$m-1$与$m+1$}{4.3(b)的m−1与m+1}}{数学阈值可确定为$m+1$；叠压字形仍标为未彻底消歧}
两个扫描的这一式均为手写改动，符号笔画叠压；英文第45页明确保留“$m-1$ or $m+1$ (illegible, ask AG)”。本次不以Dot118所称“法68较清晰”替代字形证据。正文采用$m+1$的理由是被引条目的数学内容：法第19页4.3(b)在(v)、(vi)的原下界为2；当$m=1$时，推广须退回此条件。同页第67/68页3.1与3.2已经明确给出$m+1$。第4节的可构造性讨论把这些一般纤维结论传到稠密开参数集，并不把下界降低为$m-1$。

下面的反例直接排除$m-1$作为所指普遍充分条件。取特征零域$k$，令
\[
 X=V(YZ-X^2)\subset\mathbf P^2_k,
 \qquad f:X\hookrightarrow\mathbf P^2_k,
 \qquad m=1.
\]
这是一条光滑、几何整的二次曲线，$f$为闭浸入，因而无分歧；像维数为1，满足错误下界$m-1=0$。参数仿射开片中的一般直线可写成$Y=aX+bZ$。它不含曲线的无穷远点，交集在$Z=1$上由
\[
 u^2-au-b=0
\]
给出。判别式$a^2+4b$在$k(a,b)$中非零，所以该截面在代数闭包上是两个互异点，不几何不可约，也不几何整。这个反例不依赖只取一个特殊坏直线。它同时说明$m$也不能在这些条件下取代$m+1$。

将同页3.1的超平面结论迭代$m$次时，经过$j$次一般截取后像维数为$d-j$；最后一步仍须有$d-(m-1)\geq2$，即$d\geq m+1$。这只是本次阈值核验与原作者给出的归约路线，不意味着3.1.1的原稿缺证或4.3的全部依赖已被重新证明。

法第18页4.2的旧(vi)已划销；其现行(v)是模的可分性，没有这里所说的两项几何整／几何不可约条件。因此原稿合写“4.2和4.3”之后的维数替换，明确对应4.3(b)(v)、(vi)。4.4(b)应逐个分支修改像维数条件，不能把连通性指标与线性截面的余维一并改名。法第20页作者对4.4补加的拟有限限制及一般态射下改看像维数的版本，仍须与其本条一起阅读。

\audit{V20-EXT-A03}{点闭包判据与独立一般超平面}{给出限定于域上射影空间的译校核验，不扩张任意集合的量词}
以下为译校论证。设$P=\mathbf P^N_k$，$0\leq m\leq N$，$B$为余维$m$线性子簇的Grassmann参数空间。对$x\in P$，令$W=\overline{\{x\}}$并赋予约化结构，$d=\dim W$。$W$上的关联空间$I_W$是相对维数$\dim B-m$的Grassmann丛，且为整概形。因此
\[
 \dim I_W=d+\dim B-m.
\]
若$d<m$，其在$B$中的像不含一般点。若$d\geq m$，射影相交维数定理保证$W$在参数函数域上的基变换与一般线性子簇相交，故$I_W\to B$占优。$I_W$的一般点同时映到$W$的一般点$x$和$B$的一般点，所以$x$本身有点在一般截面中映到它。于是
\[
 x\text{在一般截面的投影像中}
 \quad\Longleftrightarrow\quad
 \dim\overline{\{x\}}\geq m.
\]
此处只用到射影相交定理的非空性部分，可核\href{https://stacks.math.columbia.edu/tag/0B2N}{The Stacks Project，33.34.3，Tag 0B2N}。

若$Z$可构造，它是有限个不可约局部闭集的并，每一项含其闭包的一般点；逐点不等式便等价于$\dim\overline Z<m$。这时$\overline Z$上的关联空间到$B$的像为真闭集，故存在一个非空开参数集避开$Z$。对任意$Z$不能这样改量词：取代数闭域上$\mathbf P^2$的全体闭点与$m=1$，每个点闭包的维数都是0，一般直线在这一集合的逆像为空；但每个闭参数的直线都含闭点，故不存在一个非空开参数集的所有截面都避开$Z$。

取独立一般超平面时，应在$(P^\vee)^m$中线性独立的有序元组开集上取一般点。把元组送到其交的映射到$B$占优，因为余维$m$子空间的消去空间总可选基。每一步的“一般”是在此前参数函数域上继续选取，不能把若干各自满足某性质的任意参数直接组合。几何不可约性等性质在这里经相应域扩张比较；两个普遍族并非因此成为同一个参数空间。关于$\dim X-m$，仅对其像维数至少$m$、从而未被一般截面避开的不可约分支使用相应维数计算；混合分支不能盲目沿用全体$X$的最大维数。

\audit{V20-EXT-A04}{换基、Koszul正则性与纤维条件}{补证明确限于适用条件；不以复形同构替代全部正则性假设}
设$(R,\mathfrak m)$为诺特局部环，$M$为有限非零$R$-模，并在截面零点处写其分量为$f_1,\ldots,f_m\in\mathfrak m$。可逆换基给出两个Koszul复形的同构：一次项使用相应的可逆矩阵，高次项使用其外幂，张量$M$后仍为同构。在这些假设下，正次数Koszul同调消失等价于这组分量为$M$-正则序列；可核\href{https://stacks.math.columbia.edu/tag/09CC}{The Stacks Project，15.31.7，Tag 09CC}。所以这一局部意义的正则性与基无关。

Dot118把上述等价标为Tag 062I；实际核阅该标签所得是“$H_1$-正则推出准正则”，不是所需的四种正则性等价，故本版改引09CC。对于任意环、任意模，换基的复形同构本身仍成立，但上述反向等价不能因此无条件成立；\href{https://stacks.math.columbia.edu/tag/062D}{Tag 062D所列的一般单向蕴涵及条件}明确区分这一点。

零点集外还须明确术语约定。例如在非零域$k$上，$(0,1)$与$(1,0)$由换基互换，两个Koszul复形均可缩。但若以“依次乘法单射、允许最终商为零”定义正则，$(1,0)$满足而$(0,1)$不满足；若采用严格要求最终商非零的正则序列定义，则二者都被排除，又不能直接等同于5.2的秩一乘法单射定义。由此不能从原稿一句“验证与基无关”就宣称任意模、任意点上的各种定义已经等价。本版保留原作者的待验提示，以上仅说明有条件的补证及所需约定。

\subsec{\texorpdfstring{$(S_m)$下的局部支集维数判据}{(Sm)下的局部支集维数判据}}
仍在上述诺特局部有限模框架内，令$I=(f_1,\ldots,f_m)\subset\mathfrak m$，并假设$M$满足$(S_m)$。则
\[
 \begin{gathered}
 f_1,\ldots,f_m\text{为}M\text{-正则序列}
 \quad\Longleftrightarrow\\
 \dim\Supp(M_{\mathfrak p})\geq m
 \quad\text{对所有}\quad
 \mathfrak p\in\Supp M\cap V(I).
 \end{gathered}
\]
这里是每个包含$I$的支集点处的局部支集维数，不能用一次全空间维数比较代替。证明如下：由$(S_m)$的定义及深度不超过局部支集维数，右侧条件等价于所有这些$M_{\mathfrak p}$的深度至少为$m$。理想深度为这些深度的下确界：取$I$内一条极大$M$-正则序列，商模上$I$不能再含非零因子；素避开给出一个包含$I$的关联素理想，在该点局部化后商模深度为0，得到所需等号。任一正则序列在上述各点局部化后仍正则，给出另一向不等式。深度定义、有限性与深度的维数界见\href{https://stacks.math.columbia.edu/tag/00LE}{Tag 00LE，10.72}；$(S_m)$的模版本定义见\href{https://stacks.math.columbia.edu/tag/031O}{Tag 031O，10.157.1}。

还需说明，为何$I$-深度至少$m$控制的是所给的$m$个生成元。用Koszul对偶复形$\Hom_R(K_\bullet(f),M)$。$I$中每个元的乘法都零伦。若$I$中有一条长度$m$的$M$-正则序列，对其首项的乘法短正合列取Koszul上同调，再逐项归纳，可得该对偶复形在次数$q<m$的上同调均为零。Koszul自对偶性把这转成$K_\bullet(f)\otimes_R M$的全部正次数同调消失，继而由上述09CC得到所给序列正则。反向则直接由这组长度$m$的正则生成元给出。这里没有把“理想内存在正则序列”偷换成“理想内任取$m$个元都正则”；所给$m$个元确实须生成$I$。

在5.5的纤维语境中，应先取$M$为$F_{k(\xi)}$在$X_{k(\xi)}$上相应的有限局部模。由$F_s$的$(S_m)$过渡到该纤维模，可用正式EGA IV第二部分6.7.1(ii)：已实际核阅其印刷第145--146页，命题要求凝聚模、局部诺特的域上预概形，且原空间局部有限型或扩域有限型。第5节的有限表示框架在域纤维上满足相应有限性。该结果给出的域扩张稳定性不能当作任意总空间、任意模上的无条件结论。

两个具体模型说明不可删除的条件。第一，令$R=k[x,y]_{(x,y)}$、$M=R$，它为Cohen--Macaulay，因而满足$(S_2)$；但$(x,xy)$不是正则序列。因为模去$x$后第二项为零，且
\[
 H_1(K_\bullet(x,xy))\simeq R/(x)\ne0,
\]
其循环由$(-y,1)$生成，边界为$x(-y,1)$的倍数。相应支集点$(x)$处维数为1，正好不满足所需的下界2。第二，令$S=\Spec k[t]$，$T=\Spec k[t,x]$，$F=\mathcal O_T$，$\phi=tx$。总空间中$tx$是非零因子，但在$t=0$的纤维上截面恒为零；余核$k[t,x]/(tx)$有非零的$t$-扭元素$[x]$，不在$S$上平坦。因此总空间正则性不能替代5.3--5.7所用的纤维正则性。这些是译校反例，不是原作者所写的证明。

\audit{V20-EXT-A05}{实际代数核验及其限度}{有限模型与符号恒等式已重跑；不据此认证全部几何推广}
本次实际重跑Dot118的\texttt{verify\_local.py}，不是沿用其历史PASS记录。程序在$q=2,3$、向量空间维数$r=2,3,4$中枚举子空间，比较线性方程核的交与矩阵秩，核验了50534对子空间；并在$r\leq3$时核验独立有序元组的像遍历相应Grassmann有限点集。这是线性相交的有限模型，不能独立证明任意域上的一般点论证或非线性簇的Bertini结论。

同次运行核验了符号恒等式：若$f=(x,y)$、$A=\left(\begin{smallmatrix}a&b\\c&d\end{smallmatrix}\right)$、$g=fA$，且$d_2(f)=(-y,x)^t$，则
\[
 A\,d_2(g)=\det(A)\,d_2(f).
\]
只有在$A$可逆时这才是换基同构；恒等式本身不证明正则性。另以独立的符号计算复核二次曲线的判别式$a^2+4b$、上面的非正则循环以及$k[t,x]/(tx)$中$t[x]=0$但$[x]\ne0$。核验对象、命令、结果和来源散列记于\texttt{recollation/reports/v20\_extensions\_findings.json}。本模块没有宣称已完成2.1--5.7全部推广的数学证明。

\subsec{保留的待核事项}
4.3的手改阈值字形仍标注叠压，虽然所指充分条件已可数学地采用$m+1$。5.2的数处划销词组与5.5被替换的旧缩写未强行复原。原作者建议明确写出的3.1.1推广、任意环境中的正则截面定义及其基无关性，仍按其具体范围分别处理；本报告的局部补证不填平原稿所有条件缺口。所有“照原样”“证明基本不变”均保持为作者的推广指示，所援原条目的假设、更正、外部引文与既有证明状态不因这句话而改变。

\subsec{旧脚注逐字存档}
以下为第十九版D17-M013至D17-M024的原文，作为历史校注存档。其“底本”主要指英文编辑稿，故有关法稿是否取得、字形或错误归属的历史状态，以本版上列实核结果为准。这里的逐字保存不表示重新认定旧注的全部判断。

\par\medskip\noindent\textbf{\texttt{D17-M013}（原注全文）。}底本简写为“$\dim x\leq m-1$”。前文2.2讨论 $P$ 中子集在一般超平面中的逆像何时为空，其 $m=1$ 情形要求每一点都闭；证明也明确取点的闭包。因此这里把点的维数解释为其在 $P$ 中闭包的维数，而不是局部环维数或余维。对于域上有限型射影空间，这也等于该点剩余域在底域上的超越次数。此为记号澄清，不是在此补造全部推广证明。

\par\medskip\noindent\textbf{\texttt{D17-M014}（原注全文）。}保留这句编写者之间的插语及原英译者对语种、语气的疑问；它不构成数学论证步骤。

\par\medskip\noindent\textbf{\texttt{D17-M015}（原注全文）。}这是一条对原命题2.6的编写指示，必须连同该命题原有适用条件以及本节相应的非空性条件阅读。不能抽离条件后声称任意 $X$ 的任意余维 $m$ 截面都具有维数 $\dim X-m$。底本第8页2.6本身还有标为字迹不清的附语，本译稿没有将其恢复。

\par\medskip\noindent\textbf{\texttt{D17-M016}（原注全文）。}2.7原文实际用 $Z_i$ 表示关联素循环，条件写为 $\dim f(Z_i)>0$；“同样的替换”是把这里的下界改为 $m$，不是将 $Z_i$ 改名为 $X_i$。本段各“有限”替换均针对相应原条目中的像集或例外集，仍须保留原条目的构造性、有限型等假设。

\par\medskip\noindent\textbf{\texttt{D17-M017}（原注全文）。}本段逐项保留原稿对2.8—2.18的指示，并未将各条的假设删去或重新统一。例如2.12原文有 $f$ 无分歧的假设。“证明基本不变”是原稿判断，不是本译稿对全部推广证明完成独立验证的声明。

\par\medskip\noindent\textbf{\texttt{D17-M018}（原注全文）。}这里保留原稿的省略号及编写建议，没有凭引理编号补造3.1.1的推广陈述。此处 $G$ 是3.3中定义的对象，不应与本节开篇的泛商模 $G$ 混同。

\par\medskip\noindent\textbf{\texttt{D17-M019}（原注全文）。}底本明确留下“$m-1$ or $m+1$ (illegible, ask AG)”的两种读法。本译稿未取得可消除歧义的原手稿，因此不选择其中一个作为正确正文条件，也不凭邻近条目的 $m+1$ 猜补；这一项仍待核。

\par\medskip\noindent\textbf{\texttt{D17-M020}（原注全文）。}底本作“[Le laius 5.1 Fr or Latin 67]”，保留着英译编辑者对语言及页码的标记；这里按其语境译作“那段论述”，不把67当作数学条目号。

\par\medskip\noindent\textbf{\texttt{D17-M021}（原注全文）。}底本的$T$排作“$X_Sx\mathrm{Gr}^m$”，此处展开为纤维积；截面所在模排作$\mathcal O_{X_{k(\xi)}}(1)_{k(\xi)}\otimes G(\xi)$，重复的基变换下标在正文整理为通常写法。$G$仍指本节开头的泛秩$m$商模，$G(\xi)$是其在$\xi$处的纤维；并非第5节中作为限制模所用的同名字母。

\par\medskip\noindent\textbf{\texttt{D17-M022}（原注全文）。}底本前半句写“rank $n$”，后半句却写“$m$ sections”；本段及相邻两段以$m$表示截面的秩。秩$m$的局部自由模在一组基下恰有$m$个分量，故正文将前一处$n$统一为$m$，而非改变所需序列的长度。底本还留有“this means that such a section if $F$ regular”的句法残损，按后文定义性语境译出。基无关性仍按原稿保留为待验证事项，不在这里宣称已给出证明。

\par\medskip\noindent\textbf{\texttt{D17-M023}（原注全文）。}底本此处排作$F_S$。所指5.5先固定$\xi$在$s\in S$之上，再假设纤维模$F_s$满足$(S_1)$；此处推广的是同一个纤维上的条件，故正文写$F_s$。原稿没有在本段给出具体维数判据，本译稿不凭概括性指示补造公式或删减前文假设。

\par\medskip\noindent\textbf{\texttt{D17-M024}（原注全文）。}“照原样”是原作者的推广指示，并不意味着可以丢弃5.6中的有限表示、相对于$S$平坦等假设或5.7中的固有性条件；本译稿仅译出指示本身，未重写这些条目的完整陈述和证明。
